%% PAPER 2 — DRAFT v0.4 (corrected)
%% A Hamiltonian Dispersion Framework for Kerr Photon Rings,
%% Frequency-Dependent Constraints, Superradiance,
%% and Eikonal Quasinormal Modes
%%
%% Companion to: CQG-117140 (Paper 1)
%% Status: DRAFT — not for distribution
%%
%% v0.4 changelog (internal review, pre-submission):
%%  - FIX 1: corrected ZAMO 4-velocity formula (Section 2.1)
%%  - FIX 2: added dimensionless Lambda-hat definition (Section 4)
%%  - FIX 3: softened abstract "constraining" claim
%%  - FIX 4: renamed Table 2 from "bound" to "sensitivity scale"
%%  - FIX 5: sharpened A3.3 vs A3.4 separation
%%  - FIX 6: softened eikonal-QNM "survives" language
%%  - FIX 7: disambiguated "validated against GR" -> Lambda->0 limit
%%  - FIX 8: superradiance qualified as kinematic-threshold-only in abstract
%%  - FIX 9: replaced literature TODO with citable paragraph
%%  - FIX 10: added Jacobian-based explanation of C_pro/C_ret asymmetry
%%  - Retained: Appendix A (Carter non-separability) from v0.3

\documentclass[12pt,a4paper]{article}
\usepackage{amsmath,amssymb,amsthm}
\usepackage{graphicx,hyperref,geometry,natbib,booktabs,xcolor}
\geometry{margin=2.5cm}

\newtheorem{theorem}{Theorem}
\newtheorem{result}{Result}
\newtheorem{remark}{Remark}
\newcommand{\TODO}[1]{\textcolor{red}{\textbf{[TODO: #1]}}}
\newcommand{\Lam}{\Lambda}
\newcommand{\kloc}{k_{\rm loc}}
\newcommand{\OmH}{\Omega_H}
\newcommand{\Omph}{\Omega_{\rm ph}}
\newcommand{\lLyap}{\lambda_{\rm Lyap}}
\newcommand{\bph}{b_{\rm ph}}

\begin{document}

\title{%
  A Hamiltonian Dispersion Framework for Kerr Photon Rings,\\
  Frequency-Dependent Shadow Constraints, Superradiance,\\
  and Eikonal Quasinormal Modes
}
\author{%
  Dimitar Kretski\\
  Center for Hydro- and Aerodynamics,\\
  Bulgarian Academy of Sciences, Varna, Bulgaria\\
  \texttt{ORCID: 0000-0001-5108-2243}
}
\date{\today\ --- \textbf{DRAFT v0.4}}
\maketitle

%% ── Abstract ────────────────────────────────────────────────────────────────
\begin{abstract}
We develop a Hamiltonian dispersion framework for massless wave
propagation in Kerr spacetime with a quartic local-frame correction,
\begin{equation*}
  H = \tfrac{1}{2}\bigl[g^{\mu\nu}k_\mu k_\nu
    + \Lambda\,(k_{\rm loc}^{\rm ZAMO})^4\bigr],
\end{equation*}
and derive four interconnected results.
\textbf{(1) Photon-ring correction:}
the impact parameter shifts as $\delta b = C(a^*,\,{\rm ring})\,\Lambda E^2$
with $C_{\rm pro}=-0.6230\,M^{-1}$ and $C_{\rm ret}=+11.495\,M^{-1}$
for $a^*=0.9$; the $E^2$ scaling is proved analytically, and the
large $|C_{\rm ret}/C_{\rm pro}|\approx18.5$ asymmetry is traced to
near-degeneracy of the ring-defining Jacobian at the retrograde ring
(Section~\ref{sec:photon_ring}), not to the local dispersive momentum
$k_{\rm loc}^4$ itself.
\textbf{(2) Chromatic shadow:}
the shadow size is frequency-dependent, with
$\delta\theta\propto\Lambda\,\omega^2/D$,
and the frequency-dependent sensitivity scale satisfies
$\Lambda_{\max}\propto\omega^{-2}$ (exact algebraic scaling,
confirmed numerically to slope $-2.0000$).
\textbf{(3) Superradiance:}
we prove that the kinematic threshold $\omega=m\Omega_H$ is not
modified by the quartic dispersion; the amplification factor away
from threshold is not addressed here.
\textbf{(4) Eikonal QNMs:}
we use an eikonal continuation of the classical photon-orbit/QNM
correspondence, $\omega_{\rm QNM}\simeq\Omph(\Lambda)\,l
-i(n+\tfrac{1}{2})\lLyap(\Lambda)$, as an effective approximation,
with dispersive corrections $\delta\Omph/\Omph^{\rm GR}\simeq-0.219\,\Lambda$
and with the $\Lambda\to0$ limit validated against exact Schwarzschild
and Kerr benchmarks.
These results provide a self-consistent theoretical framework for
deriving observational tests and frequency-dependent sensitivity
estimates for the dispersion parameter $\Lambda$ through black-hole
shadow chromaticity, ring asymmetry, and gravitational-wave ringdown,
without yet constituting a measured or fitted observational bound.
\end{abstract}

\newpage\tableofcontents\newpage

%% ════════════════════════════════════════════════════════════════════════════
\section{Introduction}
%% ════════════════════════════════════════════════════════════════════════════

In a companion paper \citep{Kretski2026a} (hereafter Paper~1), we
introduced a one-parameter Hamiltonian model for dispersive massless-wave
propagation in Schwarzschild spacetime and proved two theorems on the
non-equivalence of field-theoretic and worldline descriptions under
quartic local-frame corrections.
The present paper extends the framework to Kerr spacetime and derives
four observational results.

The central Hamiltonian is
\begin{equation}\label{eq:H}
  H = \tfrac{1}{2}\bigl[g^{\mu\nu}k_\mu k_\nu
    + \Lambda\,(k_{\rm loc}^{\rm ZAMO})^4\bigr],
\end{equation}
where $k_{\rm loc}^{\rm ZAMO}$ is the locally measured wavenumber
with respect to the ZAMO observer (Section~\ref{sec:setup}),
and $\Lambda>0$ has dimensions of length squared.

We address four questions:
\begin{enumerate}
  \item[(Q1)] How does $\Lambda$ shift the photon-ring impact parameter,
    and what is the energy dependence?
  \item[(Q2)] Is the black-hole shadow chromatic?
  \item[(Q3)] Does the superradiant threshold $\omega<m\OmH$ receive
    a $\Lambda$-dependent correction?
  \item[(Q4)] What is the eikonal quasinormal-mode (QNM) prediction?
\end{enumerate}

\textbf{Summary of results.}
Q1 and Q2 are answered by a single analytical result:
$\delta b = C\,\Lambda\,E^2$ (Section~\ref{sec:photon_ring}
and \ref{sec:chromatic}).
Q3 is answered negatively by an exact proof (Section~\ref{sec:superrad}).
Q4 is answered in the eikonal ($l\gg 1$) limit
(Section~\ref{sec:qnm}).

The Event Horizon Telescope has produced resolved shadow images of
M87* \citep{EHT2019} and Sgr~A*, providing the observational context
for photon-ring-based tests of strong-field gravity. Lorentz-invariance-
violating dispersion relations have been constrained via Fermi-LAT
gamma-ray-burst timing \citep{Abdo2009} and are classified within the
Standard Model Extension framework \citep{Liberati2013}. The eikonal
correspondence between photon-ring orbital/Lyapunov parameters and
black-hole quasinormal modes was established by
\citet{Cardoso2009} and reviewed by \citet{Berti2009}; we adopt this
correspondence as the basis for Section~\ref{sec:qnm}.

%% ════════════════════════════════════════════════════════════════════════════
\section{Setup: Kerr + Quartic Dispersion}
\label{sec:setup}
%% ════════════════════════════════════════════════════════════════════════════

\subsection{Kerr metric and ZAMO observer}

We use Boyer-Lindquist coordinates with signature $(-,+,+,+)$,
$c=1$, $G=1$, and natural units $M=1$ (results are easily dimensionalised).
The spin parameter satisfies $a^*=a/M\leq 1$.

The ZAMO (zero-angular-momentum observer) has zero angular momentum,
$g_{t\phi}n^t+g_{\phi\phi}n^\phi=0$, giving angular velocity
\begin{equation}
  \Omega_{\rm ZAMO} = \frac{n^\phi}{n^t} = -\frac{g_{t\phi}}{g_{\phi\phi}},
\end{equation}
and 4-velocity
\begin{equation}
  n^\mu = \bigl(n^t,\,0,\,0,\,\Omega_{\rm ZAMO}\,n^t\bigr),
  \qquad
  n^t = \bigl(-g_{tt}-2g_{t\phi}\Omega_{\rm ZAMO}
    -g_{\phi\phi}\Omega_{\rm ZAMO}^2\bigr)^{-1/2},
\end{equation}
normalised so that $g_{\mu\nu}n^\mu n^\nu=-1$. This construction
remains valid in the ergosphere, provided $r$ is outside the event
horizon --- unlike the static observer, which ceases to exist there.
Contracting with the wave covector $k_\mu=(-\omega,\,k_r,\,k_\theta,\,m)$
gives
\begin{equation}
  n^\mu k_\mu = n^t\,\bigl(-\omega + m\,\Omega_{\rm ZAMO}\bigr).
\end{equation}
This is essential because the prograde photon ring of a rapidly
spinning Kerr black hole lies inside the ergosphere
(for $a^*=0.9$: $r_{\rm ph}^{\rm pro}=1.558\,M<r_{\rm erg}^{\rm eq}=2M$).

\subsection{Local wavenumber and Hamiltonian}

The locally measured wavenumber squared is
\begin{equation}
  (k_{\rm loc}^{\rm ZAMO})^2
  = g^{\mu\nu}k_\mu k_\nu + (n^\mu k_\mu)^2,
\end{equation}
which reduces correctly to $k^2/\alpha^2$ in the Schwarzschild equatorial
limit (Paper~1, Appendix~A).
The Hamiltonian \eqref{eq:H} satisfies $H=0$ (null constraint),
and reduces to the standard Kerr null geodesic at $\Lambda=0$.

%% ════════════════════════════════════════════════════════════════════════════
\section{Photon Ring: First-Order Correction}
\label{sec:photon_ring}
%% ════════════════════════════════════════════════════════════════════════════

\subsection{GR benchmark}

For $\Lambda=0$, the circular photon orbits satisfy $H=0$ and
$\partial H/\partial r=0$.
In the equatorial plane of Kerr with $a^*=0.9$, this gives
\citep{Bardeen1973}:
\begin{align}
  r_{\rm ph}^{\rm pro} &= 1.557855\,M, &
  b_{\rm ph}^{\rm pro} &= 2.844421\,M, \\
  r_{\rm ph}^{\rm ret} &= 3.910268\,M, &
  b_{\rm ph}^{\rm ret} &= -6.832319\,M.
\end{align}
These match the Bardeen polynomial $r^3-3Mr^2\pm 2a\sqrt{Mr^3}=0$
to residual $<3\times10^{-7}$. This reproduction of the $\Lambda\to0$
limit constitutes the GR validation of the numerical implementation
used throughout this section.

\subsection{Dispersive correction}

\begin{result}[Photon-ring shift]
\label{res:pring}
For $\Lambda>0$, the impact parameter satisfies
\begin{equation}\label{eq:db}
  \bph(E,\Lambda) = b_{\rm GR} + C(a^*,\,{\rm ring})\,\Lambda\,E^2
  + O(\Lambda^2),
\end{equation}
where $E$ is the photon energy in units $Mc^2/\hbar$, and
\begin{equation}
  C_{\rm pro} = -0.6230\;M^{-1},\qquad
  C_{\rm ret} = +11.495\;M^{-1}.
\end{equation}
The ratio $|C_{\rm ret}/C_{\rm pro}| = 18.45$ is universal for $a^*=0.9$.
\end{result}

The coefficient $C$ is obtained from implicit differentiation of the system
$\{H=0,\,\partial H/\partial r=0\}$ with respect to $\Lambda$, giving
\begin{equation}\label{eq:C_jacobian}
  C = -\frac{J_{21}\,\mathrm{rhs}_1 - J_{11}\,\mathrm{rhs}_2}{\det J},
\end{equation}
where $J$ is the Jacobian of the system and
$\mathrm{rhs}_1 = -\tfrac{1}{2}(k_{\rm loc}^2)^2$ at the GR ring.
The signs are $C_{\rm pro}<0$ (shadow shrinks) and $C_{\rm ret}>0$
(shadow grows), reflecting the opposite orientations of the two orbit families
relative to the frame-dragging.

\begin{remark}[Origin of the $C_{\rm pro}/C_{\rm ret}$ magnitude asymmetry]
\label{rem:asymmetry}
The large asymmetry between $C_{\rm pro}$ and $C_{\rm ret}$
(a factor $\approx18.5$) is \emph{not} primarily a consequence of the
local dispersive momentum $k_{\rm loc}^4$, which is in fact larger
at the retrograde ring by only a factor of $\approx2.2$
($k_{\rm loc}^2=2.827$ retrograde vs.\ $1.892$ prograde) and would,
taken alone, predict a \emph{smaller} shift there. Since
$\partial_r H=0$ identically at any photon ring
(the defining condition $J_{11}\equiv\partial_r H=0$), the
determinant in Eq.~\eqref{eq:C_jacobian} reduces to
$\det J\approx-(\partial_b H)(\partial_r^2 H)$. Direct numerical
evaluation gives $|\det J|$ smaller at the retrograde ring by a
factor of $\approx565$ relative to prograde, driven mainly by
$\partial_r^2 H$ being $\approx68$ times smaller there --- i.e., a
flatter radial effective potential at the retrograde ring. This
near-degeneracy of the ring-defining Jacobian, rather than the
magnitude of $k_{\rm loc}^4$, dominates the observed asymmetry.
Because $\partial_r^2 H$ also controls the orbital Lyapunov exponent
used in Section~\ref{sec:qnm}, this asymmetry is expected to
correlate with a correspondingly smaller $\lambda_{\rm Lyap}$ at the
retrograde ring, providing an internal consistency check between the
present section and Section~\ref{sec:qnm}.
\end{remark}

Table~\ref{tab:pring} gives numerical values for selected $\Lambda$.

\begin{table}[h]
\centering
\caption{Impact parameter $b_{\rm ph}$ vs $\Lambda$ for $a^*=0.9$, $M=1$, $E=1$.}
\label{tab:pring}
\begin{tabular}{lcccc}
\toprule
$\Lambda$ & $r_{\rm pro}$ & $b_{\rm pro}$ & $r_{\rm ret}$ & $b_{\rm ret}$ \\
\midrule
$0$        & 1.557855 & $+$2.844421 & 3.910268 & $-$6.832319 \\
$10^{-4}$  & 1.557880 & $+$2.844359 & 3.909626 & $-$6.833270 \\
$10^{-3}$  & 1.558104 & $+$2.843798 & 3.903876 & $-$6.841784 \\
$10^{-2}$  & 1.560339 & $+$2.838180 & 3.849229 & $-$6.946373 \\
\bottomrule
\end{tabular}
\end{table}

%% ════════════════════════════════════════════════════════════════════════════
\section{Chromatic Shadow}
\label{sec:chromatic}
%% ════════════════════════════════════════════════════════════════════════════

\subsection{Analytical proof of $E^2$ scaling}

\begin{theorem}[Chromatic shadow]
\label{thm:chromatic}
The first-order correction $\partial b/\partial\Lambda|_{\Lambda=0}$
is exactly proportional to $E^2$:
\begin{equation}
  \frac{\partial \bph}{\partial\Lambda}\Bigg|_{\Lambda=0} = C\,E^2.
\end{equation}
\end{theorem}

\begin{proof}
At the GR photon ring, $H_{\rm GR}=0$ implies $g^{\mu\nu}k_\mu k_\nu=0$.
Therefore
$k_{\rm loc}^2 = (n^\mu k_\mu)^2 = E^2\,K^2$, where
$K=-n^t + n^\phi b$ is an $E$-independent geometric factor.
Consequently $k_{\rm loc}^4=E^4 K^4$, all Jacobian entries scale as $E^2$,
and the implicit-differentiation formula gives
$\partial b/\partial\Lambda = C({\rm geom})\cdot E^2$ exactly.
\end{proof}

\subsection{Numerical confirmation}

The power-law fit over five decades of $E$ gives exponent $n=2.000$
for both prograde and retrograde channels.
The GR control (at $\Lambda=0$) is achromatic to
$\sigma(b_{\rm GR})<10^{-14}$ across $E\in[0.3,10]$.

\subsection{Physical formula (SI)}

To restore physical units unambiguously, define the dimensionless
dispersion parameter and photon energy
\begin{equation}
  \hat\Lambda \equiv \frac{\Lambda_{\rm SI}}{M_{\rm geom}^2},
  \qquad
  M_{\rm geom}\equiv\frac{GM}{c^2},
  \qquad
  \mathcal{E}\equiv\frac{\hbar\omega}{Mc^2},
\end{equation}
so that all coefficients $C_{\rm pro},C_{\rm ret}$ quoted above are the
dimensionless $C$ in $\delta b/M_{\rm geom}=C\,\hat\Lambda\,\mathcal{E}^2$.
Reinstating physical units:
\begin{equation}\label{eq:dtheta}
  \delta\theta_{\rm shadow} = C\,\frac{GM}{c^2}\,
  \frac{\Lambda_{\rm SI}\,(\hbar\omega)^2}{M^2c^4\,D},
\end{equation}
where $D$ is the distance to the black hole.
This formula enables a direct test: measuring the shadow size at two
radio frequencies $\omega_1,\,\omega_2$ gives
\begin{equation}
  \Lambda_{\rm SI} = \frac{\delta b(\omega_2)-\delta b(\omega_1)}
    {C\,(GM/c^2)\,(\hbar/Mc^2)^2\,(\omega_2^2-\omega_1^2)}.
\end{equation}

%% ════════════════════════════════════════════════════════════════════════════
\section{Frequency-Dependent Observational Sensitivity}
\label{sec:constraints}
%% ════════════════════════════════════════════════════════════════════════════

From \eqref{eq:dtheta} and the adopted 10\% shadow-size tolerance
(calibrated to EHT M87* uncertainty at 230\,GHz):

\textbf{Fixed-tolerance scaling (A3.3) --- primary theoretical result:}
\begin{equation}
  \Lambda_{\max}(\nu) = A\,\nu^{-2},
\end{equation}
with numerical values in Table~\ref{tab:lmax}.
The measured slope $-2.0000$ reproduces the exact $\omega^{-2}$ algebraic
scaling derived in Theorem~\ref{thm:chromatic}.

\begin{table}[h]
\centering
\caption{Illustrative fixed-tolerance sensitivity scale for $\Lambda$
(retrograde channel). These values are \emph{not} observational
constraints; they represent the value of $\Lambda$ corresponding to
the adopted 10\% angular-tolerance model, not a fit to actual EHT data
or emission physics. See text.}
\label{tab:lmax}
\begin{tabular}{lc}
\toprule
Frequency & $\Lambda_{\max}$ [m$^2$] \\
\midrule
86\,GHz  & $1.88\times10^{-8}$ \\
129\,GHz & $8.37\times10^{-9}$ \\
230\,GHz & $2.63\times10^{-9}$ \\
345\,GHz & $1.17\times10^{-9}$ \\
\bottomrule
\end{tabular}
\end{table}

\begin{remark}
Table~\ref{tab:lmax} is an internal consistency test of the model
scaling, not an independent EHT observational bound.
A rigorous constraint requires per-band sensitivity modelling,
atmospheric opacity, and imaging systematics of the EHT array.
\end{remark}

\textbf{Resolution-scaled toy extension (A3.4) --- illustrative only:}
Adopting a diffraction-limited toy model for the per-band tolerance
gives an effective slope of $-2.495$ (steeper than the algebraic $-2$).
This result depends on an assumed diffraction-limited ring-fraction
model and does not incorporate frequency-dependent instrumental
sensitivity, opacity, or calibration systematics of actual
observations. A3.3 is the theoretical scaling result of this paper;
A3.4 is retained solely as an illustrative sensitivity-modelling
extension and should not be read as a physical prediction.

\textbf{TeV feasibility (A3.3 companion test):}
A separate feasibility analysis confirms the theoretical $E^2$
chromatic-shift scaling numerically (measured slope $-2.000000$
against the expected $-2.000000$) while showing that high photon
energy alone does not provide spatial resolution: a direct
observational constraint on $\Lambda$ requires the energy-dependent
photon-ring displacement to be spatially resolved, which current
TeV instruments do not provide model-independently. Model-dependent
spectral or flux constraints are not excluded and would require a
dedicated emission and propagation model, which is beyond the scope
of this paper.

\textbf{Comparison with Fermi\,LAT:}
From Paper~1, $\Lambda\lesssim\mathcal{O}(10^{-53})$\,m$^2$
from quadratic LIV bounds \citep{Abdo2009}.
The EHT radio sensitivity scales in Table~\ref{tab:lmax} are weaker by
$\sim45$ orders of magnitude, because radio photons have
$E_{\rm photon}\ll M_{\rm BH}c^2$.

%% ════════════════════════════════════════════════════════════════════════════
\section{Superradiance}
\label{sec:superrad}
%% ════════════════════════════════════════════════════════════════════════════

\begin{theorem}[Superradiant threshold]
\label{thm:superrad}
The superradiant threshold $\omega=m\OmH$ is not modified by the
quartic dispersion $\Lambda\,\kloc^4$:
\begin{equation}
  \omega_{\rm thresh}(\Lambda) = m\OmH + O(\varepsilon_{\rm num}),
\end{equation}
where $\varepsilon_{\rm num}$ denotes the numerical proximity
$r=r_++\epsilon$.
\end{theorem}

\begin{proof}
At $r=r_+$, the ZAMO angular velocity satisfies
$\Omega_{\rm ZAMO}\to\OmH$ exactly (since $\Delta\to 0$).
Consequently, at $\omega=m\OmH$:
\begin{equation}
  n^\mu k_\mu = -n^t\,\omega + n^\phi\,m
    = n^t\,m\,(-\OmH + \Omega_{\rm ZAMO}) \to 0.
\end{equation}
Therefore $\kloc^2\to 0$, and $\Lambda\,\kloc^4\to 0$ identically at
the threshold.
The threshold condition $\omega=m\OmH$ is therefore unchanged.
\end{proof}

\begin{remark}
This result holds for the kinematic threshold only.
The rate of energy extraction (amplification factor) may receive
$\Lambda$-dependent corrections, which require computing the
full transmission amplitude of the dispersive wave equation.
\end{remark}

%% ════════════════════════════════════════════════════════════════════════════
\section{Eikonal Quasinormal Modes}
\label{sec:qnm}
%% ════════════════════════════════════════════════════════════════════════════

\subsection{Framework}

We use an eikonal continuation of the classical photon-orbit/QNM
correspondence as an effective approximation for the dispersive
Hamiltonian. In the eikonal limit $l\gg 1$, the QNM frequency is
estimated by \citep{Mashhoon1985,Cardoso2009}:
\begin{equation}\label{eq:qnm}
  \omega_{\rm QNM} \simeq \Omph(\Lambda)\,l
    - i\bigl(n+\tfrac{1}{2}\bigr)\lLyap(\Lambda),
\end{equation}
where the orbital frequency $\Omph=\dot\phi/\dot t$ and the
Lyapunov exponent
\begin{equation}
  \lLyap = \frac{\sqrt{R''(r_{\rm ph})/2}}{r_{\rm ph}^2\,|dt/d\lambda|}
\end{equation}
are now $\Lambda$-dependent through the shift in $r_{\rm ph}(\Lambda)$.
Here $R(r)$ is the standard Carter radial potential evaluated along the
$\Lambda$-shifted photon orbit. A first-principles derivation of the
dispersive QNM spectrum from the underlying radial wave equation ---
required to assess the validity of this continuation beyond the
strict eikonal regime --- is deferred to Paper~3.

\subsection{GR validation}

The $\Lambda\to0$ limit of the framework is validated against three
exact benchmarks:

\textbf{Test~1 (Schwarzschild):}
$r_{\rm ph}=3M$, $b_{\rm ph}=3\sqrt{3}\,M$,
$\Omph=\lLyap=1/(3\sqrt{3}\,M)$,
reproduced to relative error $<2\times10^{-7}$. \checkmark

\textbf{Test~2 (Kerr, $a^*=0.9$):}
$r_{\rm ph}^{\rm pro}$, $b_{\rm ph}^{\rm pro}$ match Bardeen (1973). \checkmark

\textbf{Test~3 (eikonal scaling):}
${\rm Re}(\omega_{\rm QNM})/l\to\Omph$ (constant) for $l=2,\ldots,100$. \checkmark

For $l=2$, the eikonal estimate gives
${\rm Re}(\omega)\simeq 0.703\,M^{-1}$ versus the exact Leaver value
$\simeq 0.568\,M^{-1}$ \citep{Berti2009} (24\% discrepancy, expected for $l=2$,
since Eq.~\eqref{eq:qnm} is only asymptotically valid for $l\gg1$).

\subsection{Dispersive corrections}

Table~\ref{tab:qnm} gives $\Omph(\Lambda)$ and $\lLyap(\Lambda)$
for the prograde ring.

\begin{table}[h]
\centering
\caption{Eikonal QNM parameters vs $\Lambda$ (prograde, $a^*=0.9$).}
\label{tab:qnm}
\begin{tabular}{lcccl}
\toprule
$\Lambda$ & $\Omph$ & $\lLyap$ & $\delta\Omph/\Omph^{\rm GR}$ \\
\midrule
$0$       & 0.351565 & 0.136080 & $0$         \\
$10^{-4}$ & 0.351558 & 0.136089 & $-0.0022\%$ \\
$10^{-3}$ & 0.351488 & 0.136171 & $-0.0219\%$ \\
$10^{-2}$ & 0.350800 & 0.136979 & $-0.2178\%$ \\
\bottomrule
\end{tabular}
\end{table}

The corrections are linear in $\Lambda$:
\begin{equation}
  \frac{\delta\Omph}{\Omph^{\rm GR}} \simeq -0.219\,\Lambda,\qquad
  \frac{\delta\lLyap}{\lLyap^{\rm GR}} \simeq +0.066\,\Lambda.
\end{equation}

\begin{remark}
Equation \eqref{eq:qnm} is an eikonal ($l\gg 1$) effective
approximation, not an exact result. For the dominant $l=2$
gravitational-wave mode, a full solution of the dispersive radial
wave equation is required and is deferred to Paper~3. The retrograde
ring's smaller $\lLyap$, relative to prograde, is consistent with the
Jacobian near-degeneracy identified in Remark~\ref{rem:asymmetry};
a direct numerical cross-check of this correlation (predicted ratio
$\lLyap^{\rm ret}/\lLyap^{\rm pro}\sim\sqrt{1/68}\approx0.12$) is left
for Paper~3 once the full Lyapunov calculation is available for both
ring families.
\end{remark}

%% ════════════════════════════════════════════════════════════════════════════
\section{Conclusions}
%% ════════════════════════════════════════════════════════════════════════════

We have derived four results for the quartic-dispersion Kerr model:

\begin{enumerate}
  \item \textbf{Photon-ring shift (Result~\ref{res:pring}):}
    $\delta b = C(a^*,\,{\rm ring})\,\Lambda\,E^2$, linear in $\Lambda$
    and quadratic in photon energy. For $a^*=0.9$:
    $C_{\rm pro}=-0.6230\,M^{-1}$, $C_{\rm ret}=+11.495\,M^{-1}$.
    The $\approx18.5$-fold magnitude asymmetry reflects near-degeneracy
    of the ring-defining Jacobian at the retrograde ring, not the
    magnitude of $k_{\rm loc}^4$ (Remark~\ref{rem:asymmetry}).

  \item \textbf{Chromatic shadow (Theorem~\ref{thm:chromatic}):}
    The $E^2$ scaling is exact and follows from the vanishing of
    $g^{\mu\nu}k_\mu k_\nu$ at the GR photon ring.
    The frequency-dependent sensitivity scale satisfies
    $\Lambda_{\max}\propto\omega^{-2}$.

  \item \textbf{Superradiance (Theorem~\ref{thm:superrad}):}
    The kinematic threshold $\omega=m\OmH$ is unchanged by $\Lambda$.
    The quartic term vanishes identically at the threshold because
    $\Omega_{\rm ZAMO}\to\OmH$ at the horizon. The amplification
    factor away from threshold is not addressed here.

  \item \textbf{Eikonal QNMs (Section~\ref{sec:qnm}):}
    We use the eikonal continuation
    $\omega_{\rm QNM}\simeq\Omph\,l-i(n+\tfrac{1}{2})\lLyap$
    as an effective approximation, with
    $\delta\Omph/\Omph^{\rm GR}\simeq-0.219\,\Lambda$.
    The $\Lambda\to0$ limit is validated to $<2\times10^{-7}$ against
    exact Schwarzschild and Kerr benchmarks.
\end{enumerate}

\textbf{Open questions:}
(i) Full $l=2$ QNM spectrum from the dispersive wave equation (Paper~3),
including a direct cross-check of the $\lLyap$ asymmetry against
Remark~\ref{rem:asymmetry}.
(ii) Dynamical consequences of Carter non-separability (Appendix~\ref{app:carter}).
(iii) Physical origin of $\Lambda$.

%% ── Appendix ────────────────────────────────────────────────────────────────
\appendix

\section{Carter Non-Separability: Exact Symbolic Result}
\label{app:carter}

For completeness we record the exact separability test, which motivates
item (ii) of the open questions in Section~\ref{sec:qnm}.

The standard Kerr Hamiltonian ($\Lambda=0$) is fully separable via the
Carter constant $Q_{\rm Carter}$:
\begin{equation}
  \{Q_{\rm Carter},\, H_{\rm GR}\} = 0 \quad (\Lambda = 0).
\end{equation}
This was verified to exact rational arithmetic using SymPy.

For $\Lambda>0$ and $a\neq 0$, exact symbolic computation gives:
\begin{equation}
  \{Q_{\rm Carter},\, H_\Lambda\}
  = -\frac{624358\ldots}{2809107\ldots} \neq 0,
\end{equation}
an exact rational number, confirming that the classical Carter invariant
is not conserved by the dispersive Hamiltonian.

\textbf{Interpretation.}
Non-vanishing of the Poisson bracket establishes that $Q_{\rm Carter}$
is not a constant of motion for $H_\Lambda$. It does \emph{not}
imply chaos: a new $\Lambda$-dependent invariant $Q_{\rm eff}(\Lambda)$
may exist near $Q_{\rm Carter}$.

\textbf{Dynamical tests.}
Poincar\'e sections (197 crossings per $\Lambda$ value, full 4D
Hamiltonian integration with corrected Kerr metric)
show closed $C$-shaped curves at all tested $\Lambda\in[0,10^{-4}]$.
Lyapunov exponents converge to $\lambda_{\max}\to 0$ for
$\Lambda\leq 10^{-3}\,M^2$, consistent with quasi-integrable or
regular dynamics at small $\Lambda$. Whether chaos appears at larger
$\Lambda$ requires extended integration times and is deferred to
Paper~3.

\textbf{Superradiance connection.}
The Carter non-separability is \emph{not} required for the
superradiance result in Section~\ref{sec:superrad}: that proof depends
only on the horizon geometry ($\Omega_{\rm ZAMO}\to\Omega_H$), which
holds independently of separability.

%% ── Bibliography ────────────────────────────────────────────────────────────
\bibliographystyle{plainnat}
\begin{thebibliography}{99}

\bibitem[Abdo et~al.(2009)]{Abdo2009}
Abdo, A.~A., et~al., 2009, \textit{Astrophys.\ J.}, 706, L138.

\bibitem[Bardeen(1973)]{Bardeen1973}
Bardeen, J.~M., 1973, in \textit{Black Holes}, ed.\ C.~DeWitt \&
B.~DeWitt (Gordon \& Breach), p.~215.

\bibitem[Berti et~al.(2009)]{Berti2009}
Berti, E., Cardoso, V., \& Starinets, A.~O., 2009,
\textit{Class.\ Quantum Grav.}, 26, 163001.

\bibitem[Cardoso et~al.(2009)]{Cardoso2009}
Cardoso, V., Miranda, A.~S., Berti, E., Witek, H., \& Zanchin, V.~T.,
2009, \textit{Phys.\ Rev.\ D}, 79, 064016.

\bibitem[EHT Collaboration(2019)]{EHT2019}
Event Horizon Telescope Collaboration, 2019,
\textit{Astrophys.\ J.\ Lett.}, 875, L6.

\bibitem[Kretski(2026a)]{Kretski2026a}
Kretski, D., 2026, Zenodo, \texttt{10.5281/zenodo.22018716}
(Paper~1, submitted to CQG).

\bibitem[Liberati(2013)]{Liberati2013}
Liberati, S., 2013, \textit{Class.\ Quantum Grav.}, 30, 133001.

\bibitem[Mashhoon(1985)]{Mashhoon1985}
Mashhoon, B., 1985, \textit{Phys.\ Rev.\ D}, 31, 290.

\end{thebibliography}

\end{document}
